\section*{Chapter \ref{ch:b3:extension-strike-slip} --- Extensional and Strike-Slip Tectonics}

\begin{solution}{exo:b3:extension-strike-slip:1}
The beds dip $50 - 35 = \ang{15}$, away from the faults; $\beta =
\sin\ang{50}/\sin\ang{35} = 0.766/0.574 = 1.34$.
\end{solution}

\begin{solution}{exo:b3:extension-strike-slip:2}
Dextral stepping right: basin. Dextral stepping left: ridge. Sinistral
stepping left: basin. Sinistral stepping right: ridge.
\end{solution}

\begin{solution}{exo:b3:extension-strike-slip:3}
$105/20 = \qty{5.25}{km/Myr} = \qty{5.25}{mm/yr}$, close to the
\qty{6}{mm/yr}: the rate need not have been constant over \qty{20}{Myr}.
\end{solution}

\begin{solution}{exo:b3:extension-strike-slip:4}
\qty{200}{\percent} means $\Delta L = 2L_0$: $L_0 = \qty{117.5}{km}$ and $L
= \qty{352.5}{km}$. In the north, $\Delta L = 0.5L_0$: $L_0 = \qty{470}{km}$
and $L = \qty{705}{km}$.
\end{solution}

\begin{solution}{exo:b3:extension-strike-slip:5}
$563/20 = 28$ to $563/15 = \qty{38}{mm/yr}$. The offset is a minimum, and
the drift measured by surveying across the fault zone also includes motion on
neighbouring faults of the system and elastic strain stored between
earthquakes.
\end{solution}

\begin{solution}{exo:b3:extension-strike-slip:6}
$t = 90 - \text{dip}$. At \ang{60}: $(1 + 0.6\times1.732)/(1 - 0.6\times
0.577) = 2.039/0.654 = 3.12$; at \ang{50}: $1.715/0.497 = 3.45$; at
\ang{40}: $1.504/0.285 = 5.28$. Lock-up at $\arctan0.6 = \ang{31.0}$.
\end{solution}

\begin{solution}{exo:b3:extension-strike-slip:7}
$\arctan0.4 = \ang{21.8}$ and $\arctan0.2 = \ang{11.3}$. A fault dipping
\ang{15} slips only if $\mu < \tan\ang{15} = 0.27$.
\end{solution}

\begin{solution}{exo:b3:extension-strike-slip:8}
$2.5\times4.0 = \qty{10.0}{km^2}$ and $2.5\times(4.0 - 1.5) =
\qty{6.25}{km^2}$.
\end{solution}

\begin{solution}{exo:b3:extension-strike-slip:9}
$132/1.19 = \qty{110.9}{km}$, stretched by \qty{21.1}{km}, which takes
$21.1/6 = \qty{3.5}{Myr}$.
\end{solution}

\begin{solution}{exo:b3:extension-strike-slip:10}
$3.0/\sin\ang{60} = \qty{3.464}{km}$ before and $3.0/\sin\ang{45} =
\qty{4.243}{km}$ after: \qty{0.779}{km} per block ($\beta = 1.22$). The beds
dip $60 - 45 = \ang{15}$. Each block turns as a whole: the rotation that
lowers the dip of faults dipping, say, east turns a horizontal line in the
block so that it rises to the east, and the beds dip west, opposite to the
faults.
\end{solution}

\begin{solution}{exo:b3:extension-strike-slip:11}
In a sinistral east--west zone $\sigma_1$ trends \ang{045}. R shears at
$\phi/2 = \ang{20}$ from the zone on the side of $\sigma_1$: strike \ang{070},
sinistral. R$'$ at $90 - 20 = \ang{70}$: strike \ang{020}, dextral. Tension
fractures parallel to $\sigma_1$: strike \ang{045}.
\end{solution}

\begin{solution}{exo:b3:extension-strike-slip:12}
For $\mu = 0.6$ the faults form at \ang{60.5} and lock at \ang{31.0}: $\beta =
\sin\ang{60.5}/\sin\ang{31.0} = 0.870/0.515 = 1.69$. For $\mu = 0.85$: from
\ang{65.2} to \ang{40.4}, $\beta = 0.908/0.648 = 1.40$. Faults with lower
friction rotate further before locking: weak rocks extend more before new
faults must form.
\end{solution}

\begin{solution}{pb:b3:extension-strike-slip:1}
\textbf{1.} The beds were horizontal and have turned with the blocks; the
faults have turned by the same angle, so $\theta_0 = 42 + 18 = \ang{60}$,
Anderson's dip for $\phi = \ang{30}$.

\textbf{2.} $\beta = \sin\ang{60}/\sin\ang{42} = 0.8660/0.6691 = 1.294$.

\textbf{3.} $24.0/1.294 = \qty{18.5}{km}$; the domain has been extended by
\qty{5.5}{km} (\qty{5.46}{km}).

\textbf{4.} $t = \ang{48}$: $(1 + 0.6\times0.900)/(1 - 0.6\times1.111) =
1.540/0.334 = 4.62$, against $3.12$ for new faults at \ang{60.5}: half as
much again.

\textbf{5.} At $\arctan0.6 = \ang{31.0}$; $\beta = \sin\ang{60}/\sin\ang{31}
= 1.68$ at most. Before that, new faults at about \ang{60} would cut the
blocks, and a second generation of dominoes would take over.

\textbf{6.} Every column of the hanging wall drops by the fall of the fault
beneath it over the heave; the fault is steepest near the surface and flattens
away from it, so the columns close to the fault drop most and those far away
least: the beds tilt towards the fault.

\textbf{7.} $D = A/H = 16.0/3.2 = \qty{5.0}{km}$.

\textbf{8.} $D$ is proportional to $A$: $\qty{5.0}{km} \pm \qty{0.5}{km}$.

\textbf{9.} $3.2\times(5.0 - 2.0) = \qty{9.6}{km^2}$.

\textbf{10.} The heave, \qty{3.2}{km}.

\textbf{11.} The fill thickens towards the fault, in a wedge, and its older
layers are more tilted than the younger ones: each layer was deposited while
the rollover and the space next to the fault kept growing.

\textbf{12.} $z = P/(\rho g) = \num{5.5e8}/(2800\times9.8) = \qty{20.0}{km}$.

\textbf{13.} $20.0/(24 - 16) = \qty{2.5}{km/Myr} = \qty{2.5}{mm/yr}$.

\textbf{14.} $20.0/\sin\ang{30} = \qty{40}{km}$ of slip, and $40\cos\ang{30} =
\qty{34.6}{km}$ of horizontal extension.

\textbf{15.} No: \ang{15} is below the lock-up dip of \ang{31}. Either the
fault slipped while it was steeper and was then rotated (\omterm{def:b3:extension-strike-slip:core-complex}{rolling hinge}), or
it was weak, with $\mu < 0.27$ from clays or fluids at high pressure.

\textbf{16.} By $45 - 15 = \ang{30}$: as the upper plate moved away, the
unloaded footwall rose by isostasy and bent, rotating the fault through the
hinge.

\textbf{17.} \qty{34.6}{km} against the half-graben's \qty{8.7}{km} (question
23): four times more. Core complexes lie where extension is greatest, as in
the central Basin and Range.

\textbf{18.} A sinistral fault stepping left: releasing, a basin opens.

\textbf{19.} $\qty{5}{mm/yr}\times\qty{1.2}{Myr} = \qty{6.0}{km}$.

\textbf{20.} $4.0\times6.0 = \qty{24}{km^2}$; $3.0 + 6.0 = \qty{9.0}{km}$ long;
$9.0/4.0 = 2.25$.

\textbf{21.} $132/10 = 13$ to $132/7 = 19$: the Dead Sea basin is far more
elongated, its offset many times its width.

\textbf{22.} A \omterm{def:b3:extension-strike-slip:bends}{restraining bend}: a ridge pushed up between the strands; in
section, a positive \omterm{def:b3:extension-strike-slip:flower}{flower structure}, the fault splaying upward into strands
with a reverse component, like the one 100 to \qty{300}{m} wide on the Dead
Sea Transform.

\textbf{23.} $3.2 + 5.46 = \qty{8.66}{km}$: the Merrow half-graben has been
extended by \textbf{\qty[text-series-to-math=true]{8.7}{km}}.
\end{solution}
